%==============================================================================
% Template PROYECTO DE TITULO UCM
% Autor: Marco Quiroz
%        Escuela de Ing Civil         
%        Departamento de Obras Civiles
%        Universidad Católica del Maule
%==============================================================================
\documentclass[12pt,letterpaper,fleqn]{report}
\usepackage{UCM}
  
\begin{document}
\pagestyle{empty}
\begin{center}
    \parbox[c][\textheight][t]{\textwidth}{
        \begin{center}
            \includegraphics[scale=0.2]{./figures/Escudo_UCM.png}\\
            \vspace{0.25cm}
            {\normalsize \bf UNIVERSIDAD CATÓLICA DEL MAULE}\\ \vspace{-0.15cm} 
            {\normalsize FACULTAD DE CIENCIAS DE LA INGENIERÍA}\\ \vspace{-0.15cm} 
            {\normalsize ESCUELA DE INGENIERÍA CIVIL}\\
            \vspace{1cm}
            {\normalsize \bf “PROYECTO DE TÍTULO”} \\
            \vspace{0.5cm}
            {\normalsize \bf NOMBRE (S) ALUMNO(S)} \\
            {\normalsize \bf PROFESOR GUÍA: NOMBRE PROFESOR} \\
            \vfill
            Proyecto de Título presentado en conformidad a los requisitos para obtener el título de Ingeniero Civil \\
            \vspace{0.5cm}
            TALCA, MES 202X
        \end{center}
        }
\end{center}

\newpage

\begin{center}
    \parbox[c][\textheight][t]{\textwidth}{
        \begin{center}
            \includegraphics[scale=0.2]{./figures/Escudo_UCM.png}\\
            \vspace{0.25cm}
            {\normalsize \bf UNIVERSIDAD CATÓLICA DEL MAULE}\\ \vspace{-0.15cm}
            {\normalsize FACULTAD DE CIENCIAS DE LA INGENIERÍA}\\ \vspace{-0.15cm}
            {\normalsize ESCUELA DE INGENIERÍA CIVIL}\\
            \vspace{1cm}
            Proyecto de Título presentado en conformidad a los requisitos para obtener el título de Ingeniero Civil \\
            \vspace{0.25cm}
            {\normalsize \bf “PROYECTO DE TÍTULO”} \\
            \vspace{0.5cm}
            {\normalsize \bf NOMBRE (S) ALUMNO(S)} \\
            {\normalsize \bf PROFESOR GUÍA: NOMBRE PROFESOR} \\
            \vspace{1cm}
            \begin{minipage}[t]{0.45\textwidth}
                \centering 
                {\bf COMISIÓN EXAMINADORA}\\
                \vspace{0.5cm}
                {Dr. / Mg. XX}\\ \vspace{0.75cm}
                {Dr. / Mg. XX}\\ \vspace{0.75cm}
                {Dr. / Mg. XX} \\ \vspace{0.75cm}
                {NOTA FINAL EXAMEN DE TÍTULO} \\
            \end{minipage}
            \begin{minipage}[t]{0.45\textwidth}
                \centering
                {\bf FIRMA} \\ \vspace{0.15cm}
                \wildcard{PROFESOR DE COMISIÓN EXTERNA} \\ % Please don't change this
                \wildcard{PROFESOR DE COMISIÓN EXTERNA} \\  % Add your second supervisor (TA) name here \\
                \wildcard{PROFESOR DE COMISIÓN INTERNA} \\ \vspace{-0.35cm}
                \wildcard{} \\
            \end{minipage} \\
            \vspace{1.5cm}
            TALCA, MES 20XX 
        \end{center}
    }
\end{center}

\cleardoublepage
%---------------------------------------------------------------------------
% Preamble
\pagenumbering{roman} 				% Begin Roman page numbering (i,ii,...)
%---------------------------------------------------------------------------
% Resumen
\chapter*{Resumen}
 \addcontentsline{toc}{chapter}{Resumen}
 \lipsum[1-2]

 \cleardoublepage
% Abstract
\chapter*{Abstract}
 \addcontentsline{toc}{chapter}{Abstract}
 \lipsum[3-4]

 \cleardoublepage
%---------------------------------------------------------------------------
% Agradecimientos
% Se puede remover si el estudiante no lo desea (máximo una página)
\chapter*{Agradecimientos}
\addcontentsline{toc}{chapter}{Agradecimientos}
\lipsum[5-6]

\cleardoublepage
%---------------------------------------------------------------------------
%The table of contents is automatically generated by the template
\setcounter{secnumdepth}{4}
\setcounter{tocdepth}{4}            % Sets the number of section levels in the table of contents to 4
\tableofcontents                    % Creates the table of contents
\cleardoublepage                    % Ends the current page

\listoffigures
\addcontentsline{toc}{chapter}{\listfigurename}
\cleardoublepage 
\listoftables
\addcontentsline{toc}{chapter}{\listtablename}
\cleardoublepage 
%---------------------------------------------------------------------------
% Capitulos
\pagestyle{fancy}               	% Fancy headings
\pagenumbering{arabic}				% Begin arabic page numbering (1,2,...)
\setlength{\parindent}{20pt}        % Sets default paragraph indentation to 20 pt 
% Capítulos
\input{Capitulos/Capitulo_1.tex}         

\input{Capitulos/Capitulo_2.tex} 

\input{Capitulos/Capitulo_3.tex} 

\input{Capitulos/Capitulo_4.tex} 

\input{Capitulos/Capitulo_5.tex}  

\newpage  

%---------------------------------------------------------------------------
% Bibliografia
% Referencias deben ser agregadas al archivo biblio.bib
\addcontentsline{toc}{chapter}{Bibliografía}
\bibliographystyle{apalike} % Estilo APA
\bibliography{biblio.bib}       % Carga las referencias desde el archivo biblio.bib y crea la sección de Bibliografía  
%---------------------------------------------------------------------------
% Anexos
\appendix  
\clearpage
\addappheadtotoc
\appendixpage
\renewcommand{\thechapter}{\Roman{chapter}}
\input{Capitulos/Anexo_I.tex}
\input{Capitulos/Anexo_II.tex}
%---------------------------------------------------------------------------
\end{document}
